% \input{\pSections "sec-backup"}

\section{Backup}

%%    %%    %%    %%    %%    %%    %%    %%    %%    %%
\subsection{Taylor Fit Redux}

\begin{frame}\frametitle{Taylor: Better Results Than Expected}
\center
	\href{https://homepage.physics.uiowa.edu/~ghowes/teach/phys7729/docs/Taylor50b.pdf\#page=15}{
	\begin{overpic}[ scale = 0.5 ]
		{\pLocalGraphics taylor/final-comment}
		%\put(50, 50)	{\colorbox{white}{$a+b$}}
	\end{overpic}}
	\ \\[10pt]
	``the agreement is better (than expected)'' \\[10pt]
	``a far less good agreement  would justify ... the dynamical picture''
\end{frame}
%
\begin{frame}\frametitle{Changes Since 1950}
\begin{enumerate}
	\item Quantify uncertainty
	\item Solve model directly
	\item Desktop computing
\end{enumerate}
\end{frame}

\begin{frame}\frametitle{Taylor II, Figure 1}
\center
	\href{https://homepage.physics.uiowa.edu/~ghowes/teach/phys7729/docs/Taylor50b.pdf\#page=3}{
	\begin{overpic}[ scale = 0.5 ]
		{\pLocalGraphics taylor/taylor-redux}
		\put(-45, 105)	{\colorbox{white}{$\frac{5}{2}\log R(t) = (\bl{\solnb} \pm \bl{\errb})\log t + (\solna \pm \erra)$}}
		\put(-15, 40)	{\rotatebox{90}{\colorbox{white}{$\frac{5}{2}\log R$}}}
		\put(37, -8)	{\colorbox{white}{$\log t$}}
	\end{overpic}}
\end{frame}
%
\begin{frame}\frametitle{Comparing Data, Solution Curves}
\center
	%\href{}{
	\begin{overpic}[ scale = 0.8 ]
		{\pLocalGraphics Taylor/Taylor-redux-true}
		\put(35, 52)	{\colorbox{white}{\rd{Original Method}}}
		\put(60, 40)	{\colorbox{white}{\bl{Least Squares}}}
		\put(60, 35)	{\colorbox{white}{\bl{Minimization}}}
		\put(50, 0)		{time, s}
		\put(0, 50)		{\rotatebox{90}{radius, cm}}
	\end{overpic}
\end{frame}
%
\begin{frame}\frametitle{Similarity Implies \href{https://mathworld.wolfram.com/Power.html}{Power}  Law Model}
\huge{
\begin{equation}
	\bl{R} = \bl{S(\gamma) \sqrt[5]{\frac{\rd{E}\, t^{2}}{\rho_{0}} }}
\tag{\ref{eq:one}}
\end{equation}
}
\end{frame}
%
\begin{frame}\frametitle{Similarity Implies \href{https://mathworld.wolfram.com/Power.html}{Power} Law Model}
\huge{
\begin{equation*}
	R = \mg{S(\gamma)} \sqrt[5]{\frac{\mg{E}\, t^{2}}{\mg{\rho_{0}}} } \quad \Rightarrow \quad
	R = \bl{a} t^{\bl{b}} 
\end{equation*}
}
\end{frame}
%
\begin{frame}\frametitle{Similarity Implies \href{https://mathworld.wolfram.com/Power.html}{\yl{Power}}  \yl{Law} Model}
% Taylor's premise
\huge{
\begin{equation}
	R(\bl{a}, \bl{b},t) = \bl{a} t^\bl{b}
\label{eq:model}
\end{equation}
}
\end{frame}
%
\begin{frame}\frametitle{Least Squares \rd{Definition}}
\huge{
\begin{equation}
	r(a, b) = \bl{a} t^{\bl{b}} - R
\label{eq:residual}
\end{equation}
\pause
\begin{equation}
	(\bl{a}, \bl{b})_{LS} \rd{\, \equiv \,} \argmin_{(\bl{a},\bl{b})\in\mathbb{R}^{2}} \normt{\bl{a} t^{\bl{b}} - R}
\label{eq:merit}
\end{equation}
}
\end{frame}
%
\begin{frame}\frametitle{Solution Using \href{https://www.wolfram.com/mathematica/	}{Mathematica}}
\center
	%\href{}{
	\begin{overpic}[ scale = 0.375 ]
		{\pLocalGraphics Taylor/mm-screen-02}
		%\put(35, 52)	{\colorbox{white}{\rd{Faux Linearization}}}
	\end{overpic}
\end{frame}
%
\begin{frame}\frametitle{Comparing Solution Points}
\center
	%\href{}{
	\begin{overpic}[ scale = 0.6 ]
		{\pLocalGraphics Taylor/Taylor-redux-both}
		\put(7, 87)		{\colorbox{white}{$\normt{r(a,b)} = \normt{a t^{b} - R}$}}
		\put(-9, 46)	{\rotatebox{90}{\colorbox{white}{$b$}}}
		\put(49, -5)	{\colorbox{white}{$a$}}
	\end{overpic}
\end{frame}
%
\begin{frame}\frametitle{Comparing Solution Points}
\center
	%\href{}{
	\begin{overpic}[ scale = 0.6 ]
		{\pLocalGraphics Taylor/Taylor-redux-both}
		\put(7, 87)		{\colorbox{white}{$r(a,b) = \normt{a t^{b} - R}$}}
		\put(-9, 46)	{\rotatebox{90}{\colorbox{white}{$b$}}}
		\put(49, -5)	{\colorbox{white}{$a$}}
		\put(100, 47)	{\colorbox{white}{\bl{$r = 510$}}}
		\put(100, 88)	{\colorbox{white}{\rd{$r =1456$}}}
	\end{overpic}
\end{frame}
%
\begin{frame}\frametitle{\href{https://en.wikipedia.org/wiki/Errors_and_residuals}{Residual Error} Vector Elements}
\center
	%\href{}{
	\begin{overpic}[ scale = 1.0 ]
		{\pLocalGraphics Taylor/Taylor-residuals-both}
		%\put(50, 50)	{\colorbox{white}{$a+b$}}
	\end{overpic}
\end{frame}
%
\begin{frame}\frametitle{Taylor's Minimization Method}
\begin{equation}
\begin{array}{rclcr}
	\bl{y(x)} &=& (\solnb \pm \errb)\bl{x} &+& (\solna \pm \erra) \\[10pt]
	\bl{\tfrac{5}{2} \log R} &=& (\solnb \pm \errb) \bl{\log t} &+& (\solna \pm \erra)
\end{array}
\label{eq:taylor-faux}
\end{equation}
\ \\[20pt]
Critique 1: $b$ is forced to be $\frac{2}{5}$ \\[10pt]
Critique 2: \href{https://mathworld.wolfram.com/Power.html}{Power}  law is nonlinear 
%Critique 3: Nonlinear transformation of nonlinear model does not restore linearity 
\end{frame}
%
\begin{frame}\frametitle{The \yl{Horror} of Logarithmic Differences}
Linear Differences: \\
\begin{equation}
	\begin{array}{ccccc | ccccc cc}
		10 &-& 1 &= &\bl{9} &10^{9} & - & 10^{8} &= &\rd{9\times 10^{8}}
	\end{array}
\label{eq:linear-diff}
\end{equation}
$$ \bl{9} << \rd{9\times 10^{8}} $$
	\ \\
	\pause
Logarithmic Differences: \\
\begin{equation}
	\begin{array}{ccccc | ccccc cc}
		\log 10 &-& \log 1 &= &\bl{1} &\log 10^{9} & - & \log 10^{8} &= &\rd{1}
	\end{array}
\label{eq:logarithmic-diff}
\end{equation}
$$ \bl{1} \overset{!}{=} \rd{1} $$
\end{frame}
%
\begin{frame}\frametitle{Faux Linear System}
\begin{equation}
\begin{array}{ccccc}
		%
	\mathbf{A} & c &=& y \\[20pt]
		%
	\mat{cc}{ 1 & \tfrac{5}{2} \log R_{1} \\ \vdots & \vdots \\ 1 & \tfrac{5}{2} \log R_{m} } &
	\mat{c}{c_{0} \\ c_{1} } &=&
	\mat{c}{\log t_{1} \\ \vdots \\ \log t_{m}}
		%
\end{array}
\label{eq:Faux}
\end{equation}
\end{frame}
%
\begin{frame}\frametitle{Normal Equations Solution}
Define the \bl{curvature matrix}:
\begin{equation}
	\alpha = \paren{\mathbf{A}^{\mathrm{T}} \mathbf{A}}^{-1}
\label{eq:curvature}
\end{equation}
\ \\[10pt]
Amplitudes\\
\begin{equation}
	c^{*} = \alpha \mathbf{A}^{\mathrm{T}} y
\label{eq:faux-normal}
\end{equation}
\ \\[10pt]
Uncertainties\\
\begin{equation}
	\sigma = \sqrt{\frac{r_{min}\cdot r_{min}}{m-n} diag\, \alpha}
\label{eq:faux-normal}
\end{equation}
\end{frame}
%
%
\begin{frame}\frametitle{Original Taylor Solution}
\begin{equation}
	c^{*} = \mat{c}{c_{0} \\ c_{1} } = \mat{r}{\solna \\ \solnb } \pm \mat{r}{ \erra \\ \errb}
%\label{eq:}
\end{equation}
\end{frame}
%
\begin{frame}\frametitle{Comparing Solutions\jumpLittle}
\begin{table}[htp]
%\caption{default}
\begin{center}
	\begin{tabular}{rrcr}
		 %
		 $a$: & $\reala$ & $\Rightarrow$ & $10^{\frac{2}{5}c_{0}} = 59804$ \\[20pt]
		 $b$: & $\realb$ & $\Rightarrow$ & $\frac{2}{5}c_{1} = 0.4058$ \\
		 %
	\end{tabular}
\end{center}
%\label{tab:label}
\end{table}%
\end{frame}

\begin{frame}\frametitle{Taylor's Result}
\center
	\href{https://homepage.physics.uiowa.edu/~ghowes/teach/phys7729/docs/Taylor50b.pdf\#page=2}{
	\begin{overpic}[ scale = 0.55 ]
		{\pLocalGraphics taylor/fit-result}
		\put(20, 35)	{\colorbox{white}{$y = m x + b \quad \Rightarrow \quad y-mx = b$}}
		\pause
		\put(32, -5)	{\colorbox{white}{\bl{$R(t=1\,s) = 583\,m$}}}
		\pause
		\put(27, -15)	{\href{https://nuclearweaponsedproj.mit.edu/fireball-size-effects}{\colorbox{white}{\rd{$R(t=1\,s) = 571 \pm 28 \,m$}}}}
	\end{overpic}}
\end{frame}
%
%
\begin{frame}\frametitle{\href{https://math.stackexchange.com/questions/670654/least-squares-solution-for-y-axb-after-logarithmic-transformation/2241027\#2241027}{
Full Example Online}}
\center
	%\href{https://math.stackexchange.com/questions/670654/least-squares-solution-for-y-axb-after-logarithmic-transformation/2241027\#2241027}{
	\href{run:../local/articles/Mathematics-Stack-Exchange.pdf}{
	\begin{overpic}[ scale = 0.325 ]
		{\pLocalGraphics mse}
		\pause
		\put(20, 19)	{\colorbox{white}{Two \rd{Wrongs} Do Not Make a \bl{Right}}}
	\end{overpic}}
\end{frame}
%

%
\begin{frame}\frametitle{Further Examples}
\begin{itemize}
	\item \href{https://math.stackexchange.com/questions/1488747/least-square-approximation-for-exponential-functions/2163333\#2163333}{Least Square Approximation for Exponential Functions}
	\item \href{https://math.stackexchange.com/questions/1934347/substitution-yielding-linear-function-for-least-squares-fitting/2170544\#2170544}{Substitution yielding linear function for least squares fitting}
	\item \href{https://math.stackexchange.com/questions/699071/using-least-squares-method-to-determine-unknown-variables/2204288\#2204288}{Using Least Squares method to determine unknown variables}
\end{itemize}
\end{frame}
%
\begin{frame}\frametitle{\href{https://www.osti.gov/biblio/1658681}{Analysis 2022}}
\center
	\href{run:../local/articles/Myers-et-al-2020-New-Yield-Estimates-for-Nuclear-Detonations-Over-Water.pdf}{
	\begin{overpic}[ scale = 0.3 ]
		{\pLocalGraphics water-05}
		%\put(50, 50)	{\colorbox{white}{$a+b$}}
	\end{overpic}}
\end{frame}

%
\begin{frame}\frametitle{Just say \yl{No} to logarithmic fits\jumpBig}
\center
\huge{
Just say \bl{NO} to logarithmic fits.
}
\end{frame}

%%    %%    %%    %%    %%    %%    %%    %%    %%    %%
%
\begin{frame}\frametitle{\yl{Vector Tools} for Nonlinear Least Squares}
\begin{enumerate}
	\item \mfortran: \href{\urlMinpack}{minpack}, \href{https://people.math.sc.edu/Burkardt/f_src/nl2sol/nl2sol.html}{\texttt{NL2SOL}}, $\dots$
	\item \julia: \href{https://juliapackages.com/p/leastsquaresoptim}{LeastSquaresOptim.jl}
	\item \mathematica: \href{https://reference.wolfram.com/language/ref/FindMinimum.html}{\texttt{FindMinimum}}
	\item \octave: \href{https://octave.sourceforge.io/optim/function/leasqr.html}{\texttt{leasqr}}
	\item \python: \href{https://docs.scipy.org/doc/scipy/reference/generated/scipy.optimize.curve_fit.html}{SciPy}
	\item \href{https://petsc.org/release/}{\bl{PETSc}}: \href{https://petsc.org/release/manualpages/Tao/TAOBRGN/}{TAOBRGN}
\end{enumerate}
\end{frame}
%
\begin{frame}\frametitle{Nonlinear Least Squares References}
\begin{enumerate}
	\item \bl{LANL Briefing}: \href{https://mads.lanl.gov/presentations/Leif_LM_presentation_m.pdf}{Numerical Optimization using the
Levenberg-Marquardt Algorithm}\\[10pt]
	\item \bl{Duke}: \href{https://people.duke.edu/~hpgavin/ExperimentalSystems/lm.pdf}{The Levenberg-Marquardt algorithm}\\[10pt]
	\item \bl{\href{https://numerical.recipes/}{Numerical Recipes}}: \href{https://www.foo.be/docs-free/Numerical_Recipe_In_C/c15-5.pdf}{Nonlinear Models}\\[10pt]
	\item \bl{Wikipedia}: \href{https://en.wikipedia.org/wiki/Levenberg\%E2\%80\%93Marquardt_algorithm}{Levenberg–Marquardt algorithm}\\[10pt]
	\item \bl{Mathworld}: \href{https://mathworld.wolfram.com/Levenberg-MarquardtMethod.html}{Levenberg-Marquardt Method}
\end{enumerate}
\end{frame}

%%    %%    %%    %%    %%    %%    %%    %%    %%    %%
\subsection{Kamm}
%
\begin{frame}\frametitle{Evaluation of S--vN--T Blast Wave Solution}
\center
	\href{https://cococubed.com/papers/kamm_2000.pdf\#page=4}{
	\begin{overpic}[ scale = 0.45 ]
		{\pLocalGraphics kamm/kamm-01}
		%\put(85, 15)	{\colorbox{white}{\cite{bethe1947blast}}}
	\end{overpic}}
	%\ \\[10pt]
	%\cite{kamm2000evaluation}
\end{frame}
%
\begin{frame}\frametitle{Evaluation of S--vN--T Blast Wave Solution}
\center
	\href{https://cococubed.com/papers/kamm_2000.pdf\#page=6}{
	\begin{overpic}[ scale = 0.45 ]
		{\pLocalGraphics kamm/kamm-2-a}
		%\put(85, 15)	{\colorbox{white}{\cite{bethe1947blast}}}
	\end{overpic}}
	%\ \\[10pt]
	%\cite{kamm2000evaluation}
\end{frame}
%
\begin{frame}\frametitle{Evaluation of S--vN--T Blast Wave Solution}
\center
	\href{https://cococubed.com/papers/kamm_2000.pdf\#page=7}{
	\begin{overpic}[ scale = 0.45 ]
		{\pLocalGraphics kamm/kamm-2-b}
		%\put(85, 15)	{\colorbox{white}{\cite{bethe1947blast}}}
	\end{overpic}}
	%\ \\[10pt]
	%\cite{kamm2000evaluation}
\end{frame}
%
\begin{frame}\frametitle{Evaluation of S--vN--T Blast Wave Solution}
\center
	\href{https://cococubed.com/papers/kamm_2000.pdf\#page=7}{
	\begin{overpic}[ scale = 0.45 ]
		{\pLocalGraphics kamm/kamm-2-c}
		%\put(85, 15)	{\colorbox{white}{\cite{bethe1947blast}}}
	\end{overpic}}
	%\ \\[10pt]
	%\cite{kamm2000evaluation}
\end{frame}
%
\begin{frame}\frametitle{Evaluation of S--vN--T Blast Wave Solution}
\center
	\href{https://cococubed.com/papers/kamm_2000.pdf\#page=7}{
	\begin{overpic}[ scale = 0.45 ]
		{\pLocalGraphics kamm/kamm-2-d}
		%\put(85, 15)	{\colorbox{white}{\cite{bethe1947blast}}}
	\end{overpic}}
	%\ \\[10pt]
	%\cite{kamm2000evaluation}
\end{frame}
%
\begin{frame}\frametitle{Evaluation of S--vN--T Blast Wave Solution}
\center
	\href{https://cococubed.com/papers/kamm_2000.pdf\#page=8}{
	\begin{overpic}[ scale = 0.45 ]
		{\pLocalGraphics kamm/kamm-2-e}
		%\put(85, 15)	{\colorbox{white}{\cite{bethe1947blast}}}
	\end{overpic}}
	%\ \\[10pt]
	%\cite{kamm2000evaluation}
\end{frame}
%
\begin{frame}\frametitle{Evaluation of S--vN--T Blast Wave Solution}
\center
	\href{https://cococubed.com/papers/kamm_2000.pdf\#page=8}{
	\begin{overpic}[ scale = 0.45 ]
		{\pLocalGraphics kamm/kamm-2-f}
		%\put(85, 15)	{\colorbox{white}{\cite{bethe1947blast}}}
	\end{overpic}}
	%\ \\[10pt]
	%\cite{kamm2000evaluation}
\end{frame}
%
\begin{frame}\frametitle{Evaluation of S--vN--T Blast Wave Solution}
\center
	\href{https://cococubed.com/papers/kamm_2000.pdf\#page=8}{
	\begin{overpic}[ scale = 0.45 ]
		{\pLocalGraphics kamm/kamm-2-g}
		%\put(85, 15)	{\colorbox{white}{\cite{bethe1947blast}}}
	\end{overpic}}
	%\ \\[10pt]
	%\cite{kamm2000evaluation}
\end{frame}
%
\begin{frame}\frametitle{Evaluation of S--vN--T Blast Wave Solution}
\center
	\href{https://cococubed.com/papers/kamm_2000.pdf\#page=8}{
	\begin{overpic}[ scale = 0.45 ]
		{\pLocalGraphics kamm/kamm-2-h}
		%\put(85, 15)	{\colorbox{white}{\cite{bethe1947blast}}}
	\end{overpic}}
	%\ \\[10pt]
	%\cite{kamm2000evaluation}
\end{frame}
%
\begin{frame}\frametitle{Evaluation of S--vN--T Blast Wave Solution}
\center
	\href{https://cococubed.com/papers/kamm_2000.pdf\#page=9}{
	\begin{overpic}[ scale = 0.45 ]
		{\pLocalGraphics kamm/kamm-2-i}
		%\put(85, 15)	{\colorbox{white}{\cite{bethe1947blast}}}
	\end{overpic}}
	%\ \\[10pt]
	%\cite{kamm2000evaluation}
\end{frame}
%
\begin{frame}\frametitle{Evaluation of S--vN--T Blast Wave Solution}
\center
	\href{https://cococubed.com/papers/kamm_2000.pdf\#page=9}{
	\begin{overpic}[ scale = 0.45 ]
		{\pLocalGraphics kamm/kamm-2-j}
		%\put(85, 15)	{\colorbox{white}{\cite{bethe1947blast}}}
	\end{overpic}}
	%\ \\[10pt]
	%\cite{kamm2000evaluation}
\end{frame}
%
\begin{frame}\frametitle{Evaluation of S--vN--T Blast Wave Solution}
\center
	\href{https://cococubed.com/papers/kamm_2000.pdf\#page=9}{
	\begin{overpic}[ scale = 0.45 ]
		{\pLocalGraphics kamm/kamm-2-k}
		%\put(85, 15)	{\colorbox{white}{\cite{bethe1947blast}}}
	\end{overpic}}
	%\ \\[10pt]
	%\cite{kamm2000evaluation}
\end{frame}
%
\begin{frame}\frametitle{Evaluation of S--vN--T Blast Wave Solution}
\center
	\href{https://cococubed.com/papers/kamm_2000.pdf\#page=17}{
	\begin{overpic}[ scale = 0.45 ]
		{\pLocalGraphics kamm/kamm-02}
		%\put(85, 15)	{\colorbox{white}{\cite{bethe1947blast}}}
	\end{overpic}}
	%\ \\[10pt]
	%\cite{kamm2000evaluation}
\end{frame}
%
\begin{frame}\frametitle{Evaluation of S--vN--T Blast Wave Solution}
\center
	\href{https://cococubed.com/papers/kamm_2000.pdf\#page=17}{
	\begin{overpic}[ scale = 0.45 ]
		{\pLocalGraphics kamm/kamm-03}
		%\put(85, 15)	{\colorbox{white}{\cite{bethe1947blast}}}
	\end{overpic}}
	%\ \\[10pt]
	%\cite{kamm2000evaluation}
\end{frame}
%
\begin{frame}\frametitle{Evaluation of S--vN--T Blast Wave Solution}
\center
	\href{https://cococubed.com/papers/kamm_2000.pdf\#page=18}{
	\begin{overpic}[ scale = 0.45 ]
		{\pLocalGraphics kamm/kamm-04}
		%\put(85, 15)	{\colorbox{white}{\cite{bethe1947blast}}}
	\end{overpic}}
	%\ \\[10pt]
	%\cite{kamm2000evaluation}
\end{frame}
%
\begin{frame}\frametitle{Evaluation of S--vN--T Blast Wave Solution}
\center
	\href{https://cococubed.com/papers/kamm_2000.pdf\#page=18}{
	\begin{overpic}[ scale = 0.45 ]
		{\pLocalGraphics kamm/kamm-05}
		%\put(85, 15)	{\colorbox{white}{\cite{bethe1947blast}}}
	\end{overpic}}
	%\ \\[10pt]
	%\cite{kamm2000evaluation}
\end{frame}
%
\begin{frame}\frametitle{Evaluation of S--vN--T Blast Wave Solution}
\center
	\href{https://cococubed.com/papers/kamm_2000.pdf}{
	\begin{overpic}[ scale = 0.45 ]
		{\pLocalGraphics kamm/kamm-06}
		%\put(85, 15)	{\colorbox{white}{\cite{bethe1947blast}}}
	\end{overpic}}
	%\ \\[10pt]
	%\cite{kamm2000evaluation}
\end{frame}
%

%%    %%    %%    %%    %%    %%    %%    %%    %%    %%
%\subsection{Retired}
\begin{frame}\frametitle{Fireball Energy}
\center
	\href{http://homepage.physics.uiowa.edu/~ghowes/teach/phys7729/docs/Taylor50a.pdf\#page=2}{
	\begin{overpic}[ scale = 0.5 ]
		{\pLocalGraphics Taylor/Taylor-02}
		%\put(50, 50)	{\colorbox{white}{$a+b$}}
	\end{overpic}}
\end{frame}
%
\begin{frame}\frametitle{Far Field Limit}
	%
	As viewing distance increases, a concentrated, spherical explosion will resemble a point source.
	%
\end{frame}
%
\begin{frame}\frametitle{Point Source}
\center
	\href{https://doi.org/10.1063/1.1722085}{
	\begin{overpic}[ scale = 0.3 ]
		{\pLocalGraphics ss/brode-02}
		%\put(85, 5)	{\colorbox{white}{\cite{bethe1944shock}}}
	\end{overpic}}
	%\ \\[10pt]
	%\cite{brode1955numerical}
\end{frame}
%
\begin{frame}\frametitle{\WikiTvNS s}
\center
	\href{\urlWikiTvNS}{
	\begin{overpic}[ scale = 0.3 ]
		{\pLocalGraphics wiki/math}
		%\put(85, 5)	{\colorbox{white}{\cite{bethe1944shock}}}
	\end{overpic}}
	%\ \\[10pt]
	%\cite{wikiTvNS}
\end{frame}
%
\begin{frame}\frametitle{\WikiTvNS s}
\center
	\href{\urlWikiTvNS}{
	\begin{overpic}[ scale = 0.3 ]
		{\pLocalGraphics wiki/edit-01}
		%\put(85, 5)	{\colorbox{white}{\cite{bethe1944shock}}}
	\end{overpic}}
	%\ \\[10pt]
	%\cite{wikiTvNS}
\end{frame}
%
\begin{frame}\frametitle{\WikiTvNS s}
\center
	\href{\urlWikiTvNS}{
	\begin{overpic}[ scale = 0.3 ]
		{\pLocalGraphics wiki/edit-02}
		%\put(85, 5)	{\colorbox{white}{\cite{bethe1944shock}}}
	\end{overpic}}
	%\ \\[10pt]
	%\cite{wikiTvNS}
\end{frame}
%
\begin{frame}\frametitle{\WikiTvNS s}
\center
	\href{\urlWikiTvNSSS}{
	\begin{overpic}[ scale = 0.3 ]
		{\pLocalGraphics wiki/self-similar}
		%\put(85, 5)	{\colorbox{white}{\cite{bethe1944shock}}}
	\end{overpic}}
\end{frame}
%
\begin{frame}\frametitle{\WikiTvNS s}
\center
	\href{\urlWikiTvNSAss}{
	\begin{overpic}[ scale = 0.3 ]
		{\pLocalGraphics wiki/edit-03}
		%\put(85, 5)	{\colorbox{white}{\cite{bethe1944shock}}}
	\end{overpic}}
\end{frame}

%%    %%    %%    %%    %%    %%    %%    %%    %%    %%
\subsection{Navier-Stokes Equations}
%
\begin{frame}\frametitle{Backup: \href{https://en.wikipedia.org/wiki/Rankine\%E2\%80\%93Hugoniot_conditions}{Rankine-Hugoniot} Jump Conditions}
\center
	Implications:
\begin{equation}
	-m^{2} = \frac{p_{2} - p_{1}} { 1 / \rho_{2} - 1 / \rho_{1} }
\tag{\ref{eq:Michelson–Rayleigh}}
\end{equation}
\begin{equation*}
\begin{array}{ccc}
		%
	\rho_{1} u_{1}^{2} + p_{1} & = & \rho_{2} u_{2}^{2} + p_{2} \\[10pt]
		%
	m u_{1} + p_{1} & = & m u_{2} + p_{2} \\[10pt]
		%
	m \paren{ u_{1} - u_{2} }& = & p_{2} - p_{1} \\[10pt]
		%
	m & = & -\frac{p_{2} - p_{1}} {u_{2} - u_{1}}\\
		%
\end{array}
\label{eq:Michelson–Rayleigh-derive}
\end{equation*}
\end{frame}
%
\begin{frame}\frametitle{Backup: \href{https://en.wikipedia.org/wiki/Rankine\%E2\%80\%93Hugoniot_conditions}{Rankine-Hugoniot} Jump Conditions}
\center
	Implications:
\begin{equation}
	h_{2} - h_{1} = \frac{1}{2} \displaystyle\paren{ 1 / \rho_{2} - 1 / \rho_{1} } \paren{p_{2} - p_{1}}
\tag{\ref{eq:enthalpy}}
\end{equation}
\begin{equation*}
\begin{array}{ccc}
		%
	\rho_{1} u_{1}^{2} + p_{1} & = & \rho_{2} u_{2}^{2} + p_{2} \\[10pt]
		%
	m u_{1} + p_{1} & = & m u_{2} + p_{2} \\[10pt]
		%
	m \paren{ u_{1} - u_{2} }& = & p_{2} - p_{1} \\[10pt]
		%
	m & = & -\frac{p_{2} - p_{1}} {u_{2} - u_{1}}\\
		%
\end{array}
\label{eq:enthalpyh-derive}
\end{equation*}
\end{frame}

%
\begin{frame}\frametitle{Useful Web Pages}
\begin{itemize}
	\item \mg{\href{https://www.continuummechanics.org}{Continuum Mechanics}:} \href{https://www.continuummechanics.org/materialderivative.html}{Material Derivative}
	\item \mg{\href{http://www.fluiddynamics.it}{Principles of Fluid Dynamics}:} \href{http://www.fluiddynamics.it/capitoli/Tota.pdf}{The Total Derivative}
	\item \mg{\href{https://physics.stackexchange.com}{Physics StackExchange}:} \href{https://physics.stackexchange.com/questions/368399/about-the-material-derivative-of-a-fluid-particles-velocity}{material derivative of a fluid particle's velocity}
	\item \mg{\href{https://en.wikipedia.org}{Wikipedia}}: \href{https://en.wikipedia.org/wiki/Gauge_covariant_derivative}{Gauge covariant derivative}
\end{itemize}
\end{frame}
%

%%    %%    %%    %%    %%    %%    %%    %%    %%    %%
\subsection{Vector Calculus}
%
\begin{frame}\frametitle{Personae Dramatis}
	%
\center
Vector functions (Roman letters, bold)
\begin{equation}
	\mathbf{A}\colon \mathbb{R}^{d} \mapsto \mathbb{R^{d}}
\label{eq:Avector}
\end{equation}
\ \\[10 pt]
Scalar functions (Greek letters)
\begin{equation}
	\mathbf{\psi}\colon \mathbb{R}^{d} \mapsto \mathbb{R}
\label{eq:Avector}
\end{equation}
\end{frame}
%
%
\begin{frame}\frametitle{Curl Theorem}
\center
Stokes Theorem
\begin{equation}
	\oint_{\partial S} \mathbf{A}\cdot \mathbf{d}l = \int_{S} \paren{ \curl{\mathbf{A}}}  \cdot \mathbf{d}S
\label{eq:curl-theorem}
\end{equation}

\end{frame}


\endinput  %  ==  ==  ==  ==  ==  ==  ==  ==  ==
